\documentstyle[% 


righttag% 


,11pt% 


,amssymb% 


,russian%               remove in Eng. 


,izvru%                 Set izven% in Eng. 


]{amsart} 


 


% {<}


 


\newcommand{\A}{\mathrm a} 


\newcommand{\re}{\operatorname{Re}} 


\newcommand{\im}{\operatorname{Im}} 


\newcommand{\tts}[1]{} 


\numberwithin{equation}{section} 


 


%\hoffset=-15mm%                  For Eng. 


 


\setcounter{page}{12} 


\god{2003} 


\nomer{1~(488)} %                        Remove in Eng. 


%\nomer{Vol.\,47, No.\,1}%               Set in Eng. 


%\str{\pageref{begin}--\pageref{end}}%  Set in Eng. 


\udk{519.61} 


 


\author{\it Г.П.\,Астраханцев} 


% NB:   ^^ remove in Eng. 


\title{О выборе почти оптимальных параметров \\


в алгоритмах типа Эрроу--Гурвица} 


\thanks{Работа выполнена при финансовой поддержке Российского фонда 


фундаментальных исследований, грант \No\,02-01-01214.} 


 


\date{} 


% \pagestyle{myheadings}%         Set in Eng. 


\pagestyle{plain} %              Remove in Eng. 


\begin{document} 


% \thispagestyle{plain}%          Set in Eng. 


\maketitle 


\label{begin} 


 


 


\section{Введение} 


 


В данной работе рассматриваются итерационные методы типа 


Эрроу--Гурвица для решения задач о нахождении седловых точек. Такие 


задачи возникают обычно при использовании техники теории 


двойственности или смешанных методов конечных элементов. Примером 


являются системы сеточных уравнений в методе конечных элементов для 


задач типа Стокса, задач упругости, для смешанной дискретизации 


эллиптических задач. Пусть $X$ и $Y$ --- конечномерные гильбертовы 


пространства со скалярными произведениями, которые будем обозначать 


$(\cdot,\cdot)$. Рассмотрим невырожденную систему линейных уравнений вида 


\begin{equation} 


L z=F, 


\end{equation} 


где 


$$ 


L\equiv \begin{pmatrix} 


A & B^* \\ -B & 0 


\end{pmatrix} 


,\quad z\equiv\begin{pmatrix} 


x \\ y 


\end{pmatrix} 


\in X\times Y,\quad F\equiv \begin{pmatrix} 


f_1 \\ f_2 


\end{pmatrix} 


. 


$$ 


Здесь $A$ --- симметричная, положительно-определенная 


$(n_1\times n_1)$-матрица, $B$ --- 


прямоугольная матрица $n_2\times n_1$, $B^*$ --- сопряженная к $B$ 


матрица. Будем 


исследовать скорость сходимости итерационного метода типа 


Эрроу--Гурвица с тремя параметрами


\begin{gather} 


\frac{Q_A(x_{n+1}-x_n)}{\tau}=-[A x_n+B^*y_n-f_1],\\ 


% 


-\frac{\alpha_1B(x_{n+1}-x_n)}{\tau}+ 


\frac{\alpha Q_B(y_{n+1}-y_n)}{\tau}=-[-B x_n-f_2];\notag \\ 


% 


Q_A,Q_B>0,\ \ \tau,\alpha,\alpha_1>0.\notag 


\end{gather} 


Эти методы известны как предобусловленный метод Ричардсона 


($\alpha_1=0$) и 


неявный метод Удзавы ($\alpha_1=\tau$); при $\alpha_1=\tau=1$ и $Q_A=A$ 


метод (1.2) известен как 


предобусловленный метод Удзавы. Введение третьего параметра $\alpha_1$ 


вызвано желанием охватить алгоритм типа неявного метода Удзавы для 


расширенного лагранжиана с параметром. Такой алгоритм исследовался, 


например, в [1] и [2]. Спектр оператора пересчета в этих работах 


совпадает со спектром оператора пересчета $T$ в методе (1.2) (операторы 


сопряжены). Введение операторов $Q_A$ и $Q_B$ отражает либо стремление 


улучшить скорость сходимости метода (предобусловливание), либо 


наличие внутренних итераций при реализации метода Удзавы. Требование 


однозначной разрешимости (1.1) приводит к тому, что выполнены 


неравенства 


\begin{equation} 


0<\A\le\frac{(Ax,x)}{(Q_A x,x)}\le1,\quad 


0<m\le\frac{(B A^{-1}B^*y,y)}{(Q_B y,y)}\le1 


\end{equation} 


для некоторых положительных чисел $\A$, $m$. Наличие 1 в правой части 


(1.3) не 


уменьшает общности, поскольку достигается простой перенормировкой $Q_A$ и 


$Q_B$. 


 


Построению и исследованию алгоритмов решения седловых задач посвящено 


множество работ, начиная с работ Эрроу, Гурвица, Удзавы (1958 г.). 


Для примера отметим работы [3]--[5]. В настоящее время для $Q_A=A$ получен 


окончательный результат о возможностях метода (1.2) [6], даны 


оптимальные значения параметров, минимизирующих спектральный радиус 


$T$. В [6] отмечено, что при оптимальных параметрах у $T$ могут быть 


клетки Жордана второго порядка. Поэтому получение оценок скорости 


сходимости вида $\|T^n\|\le c q^n$ является сомнительной задачей, более 


вероятна 


оценка $\|T^n\|\le c n q^n$. Предлагаемая работа является попыткой 


продолжить эти 


исследования и придать характер неулучшаемости оценкам скорости 


сходимости при произвольном предобусловливателе $Q_A$ для первого блока 


уравнений. 


 


Сравнение полученных результатов будет проводиться с решением 


следующей модельной задачи. Пусть $Q_A=I$, $Q_B=I$; $A$, $B$ --- 


диагональные 


матрицы размера $n_1\times n_1$ и $n_2\times n_1$ вида ($n_2\le n_1$) 


\begin{gather} 


A= 


\begin{pmatrix} 


a_1 & 0 & 0 & \cdot & 0 & 0 \\ 


0 & a_2 & 0 & \cdot & 0 & 0 \\ 


0 & 0 & a_3 & \cdot & 0 & 0 \\ 


\hdotsfor[1.5]{6} \\ 


0 & 0 & 0 & \cdot & a_{n_1-1} & 0 \\ 


0 & 0 & 0 & \cdot & 0 & a_{n_1} 


\end{pmatrix} 


,\qquad B= 


\begin{pmatrix} 


\sqrt{b_1} & 0 & \cdot & 0 & \cdot & 0 \\ 


0 & \sqrt{b_2} & \cdot & 0 & \cdot & 0 \\ 


\hdotsfor[1.5]{6} \\ 


0 & 0 & \cdot & \sqrt{b_{n_2}} & \cdot & 0 


\end{pmatrix} 


, \\ 


% 


0<\A\le a_i\le 1;\ \ 0<m<\frac{b_i}{a_i}\equiv t_i\le1.\notag 


\end{gather} 


Для этого примера собственные значения 


оператора пересчета описываются выражением 


$$ 


\lambda=1-\tau\theta\pm\tau 


\sqrt{\theta^2-\frac{t a}{\alpha}},\quad 


\theta\equiv\frac{a}{2}(1+\alpha_1 t), 


$$ 


где $t=0$ либо $m\le t<1$. Решение задачи асимптотической оптимизации 


\begin{equation} 


q_0\equiv\min_{\tau,\alpha,\alpha_1} 


\max\begin{Sb} t\in[m,1] \\ a\in[\A,1]\end{Sb} 


\{|1-\tau a|,\ |\lambda(t,a)|\}, 


\end{equation} 


приведено в [2], $q_0(\A,m)<1$ является корнем кубического уравнения 


\begin{equation} 


q^3_0-q^2_0(1+\A)\bigg(\frac{1+m}{1-m}\bigg)+ 


q_0(2\A-1)+(1-\A)\bigg(\frac{1+m}{1-m}\bigg)=0. 


\end{equation} 


Анализ этого модельного примера приводит к выводу, что третий 


параметр $\alpha_1$ в некотором смысле лишний ($\alpha_1=\tau$), т.\,е. 


оптимальным 


методом является неявный метод Удзавы. Поведение оптимальных 


параметров $\alpha_0$, $\tau_0$ в модельном примере нельзя описать одной 


формулой 


при малых $\A$, $m$, поскольку 


\begin{alignat*}{3} 


\tau_0&\approx 2\sqrt{2}\sqrt{\frac{m}{\A}}, &\quad 


\alpha_0&=\frac{2m}{\A}, &\quad m/\A&\ll1; \\ 


% 


\tau_0&\approx2, &\quad \alpha_0&\approx 1, &\quad 


m/\A&\gg 1. 


\end{alignat*} 


Ситуация с выбором параметров аналогична методу верхней релаксации 


SOR (блочно-трех\-диа\-го\-наль\-ный случай). Имеется большой диапазон 


параметров, где качество сходимости примерно одинаковое, и 


существенно меньшая область, где происходит изменение асимптотики 


скорости сходимости. 


 


Результатом данной работы является утверждение, что в общем случае 


оценка скорости сходимости (1.2) в предположении (1.3) не намного 


хуже решения в случае модельной задачи. В [7] для смешанного метода 


конечных элементов для бигармонического уравнения дана оценка 


скорости сходимости метода типа (1.2). Было отмечено, что введение 


параметров меняет асимптотику показателя скорости сходимости при $m\ll1$. 


Используемая там техника получения результата применяется в [8] и 


данной работе. Оценка скорости сходимости алгоритма (1.2) может быть 


получена анализом резольвенты оператора $T$. Окончательная оценка 


имеет вид $\|T^n\|\le C(1-\delta)^n$, $\delta>0$, где $\delta$ 


определяется спектром, а $C$ --- нормой 


резольвенты. Будет указан такой выбор параметров $\tau$, $\alpha$, 


$\tau_1$, что при 


малых $\A$, $m$ имеют место неравенства, характеризующие качество 


итерационного процесса, 


\begin{equation} 


\begin{split} 


1-\sqrt{2\A m}&\le q_0\le(1-\delta)=1-\sqrt{2\A m}/ 


(72\sqrt{2}),\quad m/\A\le 1/4; \\ 


% 


1-\A&\le q_0\le(1-\delta)=1-\A/144, \quad \qquad\qquad m/\A\ge 1/4. 


\end{split} 


\end{equation} 


Заметим, что применяемая техника не требует симметрии оператора $A$, 


достаточно, чтобы была симметрична его основная часть. 


 


\section{Вспомогательные результаты} 


 


Оператор пересчета метода (1.2) имеет вид $T=I-\tau D^{-1}L$, где 


$$ 


D\equiv 


\begin{pmatrix} 


Q_A & 0 \\ -\alpha_1B & \alpha Q_B 


\end{pmatrix} 


. 


$$ 


Для оценки скорости сходимости будет исследован спектр оператора 


$D^{-1}L$. 


Тогда использование для представления степеней оператора $T$ 


интегральной формулы Коши 


$$ 


T^n=-\frac{1}{2\pi i}\oint_\Gamma(1-\tau\lambda)^n 


[D^{-1}L-\lambda I]^{-1}d\lambda 


$$ 


дает возможность получать оценки скорости сходимости. Далее будет 


показано, что $\Gamma$ --- это парабола, являющаяся границей области 


$\Pi$, содержащей спектр $D^{-1}L$. Для оптимального выбора параметра 


$\tau$ можно использовать решение оптимизационной задачи. Введем 


обозначения


$$


\tau_0\equiv2(\lambda^-_1+\lambda^+_1+p)^{-1}, 


\ \ \tau_-\equiv(\lambda^-_1+p)^{-1}, 


\ \ \rho(\tau,\lambda)\equiv|1-\tau\lambda|. 


$$


 


\begin{thm} 


Пусть $\Pi$ --- множество комплексных чисел 


$\lambda=\lambda_1+i\lambda_2$, 


$$ 


\Pi\equiv\{\lambda:0<\lambda^-_1\le\lambda_1\le\lambda^+_1, 


\ |\lambda_2|^2\le p\lambda_1\},\ \ p>0, 


$$ 


тогда решение минимаксной задачи 


\begin{equation} 


q_1(\lambda^-_1,\lambda^+_1,p)=\min_{\tau>0} 


\max_{\lambda\in\Pi}|1-\tau\lambda|\equiv 


\min_{0<\tau}\rho(\tau) 


\end{equation} 


имеет вид 


\begin{equation} 


\begin{split} 


q^2_1&=(1+\lambda^-_1/p)^{-1}=|1-\tau_-\lambda^-_1|^2\quad 


\text{при}\ \ p\ge\lambda^+_1-\lambda^-_1;\\ 


% 


q^2_1&=|1-\tau_0\lambda^-_1|^2 \qquad\qquad\qquad \qquad


\text{при}\ \ p\le\lambda^+_1-\lambda^-_1. 


\end{split} 


\end{equation} 


\end{thm} 


 


\begin{pf} 


Ясно, что $\rho(\tau)$ является значением $|1-\tau\lambda|$ при 


$\lambda$, лежащей на параболе. На ней 


\begin{equation} 


\rho^2(\tau,\lambda)=(1-\tau\lambda_1)^2+\tau^2p\lambda_1\equiv 


\rho^2_1(\tau,\lambda_1). 


\end{equation} 


Анализируя (2.3) для всевозможных значений $\tau$, получим 


\begin{gather} 


\rho^2(\tau)=\rho^2_1(\tau,\lambda^-_1),\quad \tau\le\tau_0;\notag\\ 


% 


\rho^2(\tau)=\rho^2_1(\tau,\lambda^+_1),\quad \tau\ge\tau_0;\\ 


% 


\rho^2(\tau)\ge1,\quad \tau\ge2/(\lambda^+_1+p). \notag


\end{gather} 


Видно, что $\tau_0$ определяется условием 


$\rho_1(\tau,\lambda^-_1)=\rho_1(\tau,\lambda^+_1)$. Поскольку 


$\tau_+=1/(\lambda^+_1+p)$, вершина 


параболы $\rho^2_1(\tau,\lambda^+_1)$ лежит левее $\tau_0$, то в (2.1) 


можно сузить область, по 


которой берется минимум $q^2_1=\min\limits_{0\le\tau\le\tau_0} 


\rho^2_1(\tau,\lambda^-_1)$, после чего, в зависимости от соотношения 


между $\tau_-$ и $\tau_0$ очевидны равенства (2.2). 


\end{pf} 


 


Получим упрощенные представления для решения модельной задачи 


$q_0(\A,m)$, 


являющиеся оценками снизу. Пусть $b\equiv(1+m)/(1-m)$, 


$\varepsilon^2_0\equiv 2\A(b-1)[3-b(1+\A)]^{-1}$. 


 


\begin{thm} 


Для решения задачи асимптотической оптимизации $(1.5)$ имеют место 


неравенства 


\begin{gather} 


(1-\A)(1+\A)\le q_0(\A,m); \\ 


% 


1-\varepsilon_0\le q_0(\A,m)\ \ \text{при} 


\ \ (1+m)/(1-m)\le(3+2\A)(1+\A)^{-2}. 


\end{gather} 


\end{thm} 


 


\begin{pf} 


Сделаем замену $q_0=1-\varepsilon$ в уравнении (1.6). Получим 


\begin{equation} 


\Phi(\varepsilon)\equiv-\varepsilon^3+\varepsilon^2 


(3-b(1+\A))+2\varepsilon(b-1)(1+\A) 


-2\A(b-1)=0. 


\end{equation} 


Пусть $\varepsilon_0$ --- корень упрощенного уравнения 


$$ 


\Phi_0(\varepsilon)\equiv\varepsilon^2(3-b(1+\A))-2\A(b-1)=0,\quad 


\Phi(\varepsilon)=\Phi_0(\varepsilon)+\Phi_1(\varepsilon), 


$$ 


a $\varepsilon_1>0$ обращает в нуль выражение $\Phi_1(\varepsilon)= 


\varepsilon[-\varepsilon^2+2(b-1)(1+\A)]$, 


$\varepsilon^2_1=2(b-1)(1+\A)$. 


Если $\varepsilon_1\ge\varepsilon_0$, то на промежутке 


$(0,\varepsilon_0)$ есть корень уравнения (2.7) (перемена знака у 


$\Phi(\varepsilon)$). Для справедливости $\varepsilon_1\ge\varepsilon_0$ 


нужно выполнение двух неравенств: $b<3/(1+\A)$, $b\le(3+2\A)(1+\A)^{-2}$, 


которые верны при условии 


\begin{equation} 


(1+m)/(1-m)\le(3+2\A)(1+\A)^{-2},\quad \A>0. 


\end{equation} 


Справедливость соотношения (2.6) установлена. Оценка (2.5) получается 


из (1.5), если положить $\alpha_1=\tau=2/(1+\A)$, $t=1$, $a=\A$, 


$\alpha=\tau/\A$. 


\end{pf} 


 


Заметим, что при $m\le1/9$ условие (2.8) становится содержательным 


($\varepsilon_0\le1$), и 


при $m$, $\A$ мaлых имеем $\varepsilon^2_0\approx2\A m$. Для всех $\A$ 


и $m\le1/9$ имеет место соотношение 


$$ 


2\A m\le\varepsilon^2_0\le9\A m. 


$$ 


 


В дальнейшем нам понадобится следствие первого из неравенств (1.3): 


$Q^{-1}_A\le A^{-1}$. Действительно, 


\begin{equation} 


(Q^{-1}_A x,x)=\sup_\varphi 


\frac{(Q^{-1}_A x,\varphi)^2}{(Q^{-1}_A\varphi,\varphi)}=\sup_u 


\frac{(x,u)^2}{(u,Q_A u)}\le\sup_u 


\frac{(A^{-1}x,x)(Ax,x)}{(Q_A u,u)}\le(A^{-1}x,x). 


\end{equation} 


 


\section{Анализ метода типа Эрроу--Гурвица} 


 


В этом пункте получим представление степеней оператора пересчета $T$. 


Для оценки $T^n$ нужно исследовать спектр оператора $D^{-1}L$. 


 


\begin{thm} 


Пусть для параметров метода $(1.2)$ верны следующие ограничения${:}$ 


$\alpha_1\le1$, 


$\alpha_1=2\sqrt{\alpha}$, тогда спектр оператора $D^{-1}L$ содержится в 


области 


\begin{equation} 


\Pi\equiv\{\lambda:0<\lambda^-_1\le\lambda_1\le\lambda^+_1, 


\ |\lambda_2|^2\le9\lambda_1/\alpha_1\}, 


\end{equation} 


где $\lambda^-_1=R(1+25/9)^{-1/2}$, $R\equiv\min\{\A/4,\ m/(\alpha4)\}$, 


$\lambda^+_1=(1+\alpha_1/\alpha)$. 


\end{thm} 


 


\begin{pf} 


1) Рассмотрим равенство 


\begin{equation} 


L z-\lambda D z=D F,\quad z,F\in X\times Y;\quad 


F\equiv\{F_1,F_2\}. 


\end{equation} 


Умножим (3.2) на $\{x,0\}$, $\{0,y\}$, $\{Q^{-1}_A B^*y,0\}$. Тогда 


получим 


\begin{gather} 


(Ax,x)+(B^*y,y)=\lambda(Q_A x,x)+(Q_A F_1,x),\\ 


% 


-(B x,y)=-\alpha_1\lambda(B x,y)+\alpha_1\lambda(Q_B y,y)- 


\alpha_1(B F_1,y)+\alpha(Q_B F_2,y), \\ 


% 


(Q^{-1}_A Ax,B^*y)+(Q^{-1}_A B^*y,B^*y)=\lambda 


(x,B^*y)+(F_1,B^*y). 


\end{gather} 


Суммируя (3.3) с комплексно сопряженным равенством (3.4) и с сопряженным 


(3.5), умноженным на $\alpha_1$, получим 


\begin{gather} 


(Ax,x)+\alpha_1(Q^{-1}_A B^*y,B^*y)=\lambda 


(Q_A x,x)+\alpha(\lambda_1-i\lambda_2) 


(Q_B y,y)-\alpha_1\ov{(Q^{-1}_A Ax,B^*y)}+\Phi(x,y),\\ 


% 


\Phi(x,y)\equiv(Q_A F_1,x)-\alpha_1\ov{(B F_1,y)}+\alpha 


\ov{(Q_B F_2,y)}+\alpha_1\ov{(F_1,B^*y)}.\notag 


\end{gather} 


 


2) Покажем, что при любых отрицательных $\lambda_1$ нет спектра. Взяв 


вещественную часть от (3.6) и используя неравенство (2.9), получим 


\begin{multline} 


(Ax,x)+\alpha_1(Q^{-1}_A B^*y,B^*y)\le \\ 


% 


\le\lambda_1(Q_A x,x)+\alpha\lambda_1(Q_B y,y)+ 


\alpha_1(Q^{-1}_A B^*y,B^*y)^{1/2}(Ax,x)^{1/2}+ 


|\re\Phi(x,y)|. 


\end{multline} 


При $\alpha_1\le4$ имеем 


$$ 


(1-\sqrt{\alpha_1}/2)[(Ax,x)+\alpha_1(Q^{-1}_A B^*y,B^*y)] 


\le|\Phi(x,y)|,\quad \lambda_1\le0. 


$$ 


 


3) Заметим, что и при больших $\lambda_1$ нет спектра. Перепишем 


вещественную часть (3.6) в виде 


$$ 


\lambda_1(Q_A x,x)+\alpha\lambda_1(Q_B y,y)=(Ax,x)+\alpha_1 


(Q^{-1}_A B^*y,B^*y)+\alpha_1\re(Q^{-1}_A Ax,B^*y)- 


\re\Phi(x,y). 


$$ 


В силу неравенств (2.9), (1.3) имеем 


\begin{multline*} 


\lambda_1(Q_A x,x)+\alpha\lambda_1(Q_B y,y)\le (Ax,x)+\alpha_1 


(A^{-1}B^*y,B^*y)+\alpha_1(A^{-1}B^*y,B^*y)^{1/2} 


(A Q^{-1}_A Ax,Q^{-1}_A Ax)^{1/2} +\\ 


% 


+|\re\Phi(x,y)|\le(Q_A x,x)+\alpha_1(Q_B y,y)+\alpha_1 


(Q_B y,y)^{1/2}(Q_A x,x)^{1/2}+|\re\Phi(x,y)|. 


\end{multline*} 


Обозначим $p^2\equiv(Q_A x,x)$, $q^2\equiv\alpha(Q_B y,y)$, 


и пусть $\mu$ --- точная постоянная в оценке 


$$ 


\frac{p^2+\alpha_1\alpha^{-1}q^2+\alpha_1/\sqrt{\alpha}\,p q} 


{p^2+q^2}\le\mu, 


$$ 


т.\,е. $\mu$ --- наибольший из $\mu_1{<}\mu_2$ корней квадратного уравнения 


$\mu^2-(1+\alpha_1/\alpha)\mu+\alpha_1(1-\alpha_1/4)/\alpha=0$. 


Предположим опять, что $\alpha_1<4$, тогда, если считать 


$\lambda_1\ge1+\alpha_1/\alpha\equiv\lambda^+_1$, 


будет иметь место неравенство 


\begin{equation} 


\varepsilon\lambda_1[(Q_A x,x)+\alpha(Q_B y,y)]\le 


|\Phi(x,y)|,\quad (1-\varepsilon)(1+\mu_1/\mu_2)=1. 


\end{equation} 


 


4) Поскольку задача (3.2) при $\lambda=0$ разрешима, покажем, что при 


малых $|\lambda|$ 


она имеет решение. Исключая $x$ из этих уравнений, получим 


\begin{gather} 


(1-\lambda\alpha_1)B(A-\lambda Q_A)^{-1}B^*y-\alpha\lambda 


Q_B y-\Phi_2=0,\quad |\lambda|<\A;\\ 


% 


\Phi_2\equiv(1-\lambda\alpha_1)B(A-\lambda Q_A)^{-1} 


Q_A F_1-\alpha_1B F_1-\alpha Q_B F_2.\notag 


\end{gather} 


С учетом разложения в ряд Неймана, получим 


\begin{multline*} 


|\big((A-\lambda Q_A)^{-1}B^*y,B^*y\big)|=\bigg| 


(A^{-1}B^*y,B^*y)+\sum^\infty_{n=1}\lambda^n 


(A^{-1}K^n B^*y,B^*y)\bigg|\ge \\ 


% 


\ge\bigg(1-\sum^\infty_{n=1}\bigg|\frac{\lambda}{\A}\bigg|^n\bigg) 


(A^{-1}B^*y,B^*y)=\bigg(1-\frac{|\lambda|/\A}{1-|\lambda|/\A}\bigg) 


(A^{-1}B^*y,B^*y), 


\end{multline*} 


где $K\equiv Q_A A^{-1}$. 


При получении последнего неравенства использовано следствие из


первого неравенства (1.3): 


$\frac{[K x,x]}{[x,x]}\equiv \frac{(A^{-1}Q_A A^{-1}x,x)} 


{(A^{-1}x,x)}\le\frac{1}{\A}$. 


Умножив (3.9) скалярно на $y$, имеем 


$$ 


|(\Phi_2,y)|\ge m|1-\lambda \alpha_1|\bigg(1- 


\frac{|\lambda|/\A}{1-|\lambda|/\A}\bigg) 


(Q_B y,y)-\alpha|\lambda|(Q_B y,y). 


$$ 


Выберем $|\lambda|$ настолько малым, чтобы выполнялось неравенство 


$m(1-\alpha_1|\lambda|)(2-1/(1-|\lambda|/\A))- 


\alpha|\lambda|\ge\varepsilon>0$. 


Для этого, например, потребуем выполнения следующих условий: 


\begin{equation} 


\alpha_1\le 1,\ \ |\lambda|\le\min\{\A/4,\ m/(4\alpha)\}\equiv R, 


\ \ \varepsilon=m/4. 


\end{equation} 


Тогда 


\begin{equation} 


|(\Phi_2,y)|\ge(Q_B y,y)m/4. 


\end{equation} 


Поскольку еще из первого блока уравнений (3.2) следует 


\begin{equation} 


(Ax,x)(1-1/4)\le(Q_A F_1,x)+(Ax,x)^{1/2} 


(Q_B y,y)^{1/2}, 


\end{equation} 


то в круге $|\lambda|\le R$ нет спектра. 


 


5) Дадим некоторое уточнение п.\,4. При малых $\lambda_1\ge0$ и вне угла, 


близкого к $\pi/4$, нет спектра. Из мнимой части равенства (3.6) можно 


найти $(Q_B y,y)$ и продолжить неравенство (3.7) 


\begin{multline} 


(Ax,x)+\alpha_1(Q^{-1}_A B^*y,B^*y)\le 2\lambda_1(Q_A x,x)+ 


\alpha_1(1+\lambda_1/|\lambda_2|)(Q^{-1}_A B^*y,B^*y)^{1/2} 


(Ax,x)^{1/2} +\\ 


% 


+|\re\Phi(x,y)|+(\lambda_1/|\lambda_2|) 


|\im\Phi(x,y)|. 


\end{multline} 


Предположим, что имеет место неравенство 


\begin{equation} 


\lambda_1/\A+\sqrt{(\lambda_1/\A)^2+\alpha_1\left(1+ 


\frac{\lambda_1}{|\lambda_2|}\right)^2\,\Big/4}\le1-\varepsilon_1, 


\end{equation} 


тогда 


\begin{equation} 


\varepsilon_1[(Ax,x)+\alpha_1(Q^{-1}_A B^*y,B^*y)]\le 


(1+\lambda_1/|\lambda_2|)|\Phi(x,y)|. 


\end{equation} 


Если считать, что справедливы соотношения 


\begin{equation} 


\alpha_1\le1,\ \ \lambda_1/|\lambda_2|\le0.6, 


\ \ 0\le\lambda_1\le\frac{\A}{4\sqrt{1+25/9}}, 


\end{equation} 


то (3.14) имеет место с $\varepsilon_1=0.06$. 


 


6) Справедливость утверждений из пп.\,4, 5 приводит к тому, что при 


малых $\lambda_1\ge0$ нет спектра. Действительно, пересечение прямой 


$|\lambda_2|=\lambda_1/0.6$ из условия 


(3.16) с окружностью $|\lambda_2|=R$ приводит к точке 


$\lambda^-_1\equiv R/\sqrt{1+25/9}$, и, очевидно, (см. (3.10)) 


$$ 


\lambda^-_1\le\frac{\A}{4\sqrt{1+25/9}}. 


$$ 


Таким образом, объединяя п.\,4--6, получим, что в полосе


$$


0\le 


\lambda_1\le \lambda^-_1= (1+25/9)^{-1/2} \min\{\A/4,\ m/(4\alpha)\}


$$ 


нет спектра ($\alpha_1\le1$). 


 


7) Покажем, что вне параболы вида $|\lambda_2|^2\le p\lambda_1$ нет 


спектра. Пусть $0<\lambda_1\le\lambda^+_1$ и более точно 


проведем оценки в (3.13). Из мнимой части равенства (3.3) найдем оценку 


$(Q_A x,x)$


$$ 


|\lambda_2|(Q_A x,x)\le (Q_A x,x)^{1/2}[(Q^{-1}_A B^*y,B^*y)^{1/2} 


+(Q_A F_1,F_1)^{1/2}]. 


$$ 


Продолжая неравенство (3.13), имеем 


\begin{gather*} 


(Ax,x)+\alpha_1(Q^{-1}_A B^*y,B^*y)\le 


(2\lambda_1/\lambda^2_2)(Q^{-1}_A B^*y,B^*y)+ \\ +\alpha_1 


(1+\lambda_1/|\lambda_2|)(Q^{-1}_A B^*y,B^*y)^{1/2} 


(Ax,x)^{1/2}+\Phi_3, \\ 


% 


\Phi_3\equiv(1+\lambda_1/|\lambda_2|)|\Phi(x,y)|+ 


(2\lambda_1/\lambda^2_2)[2(Q^{-1}_A B^*y,B^*y)^{1/2} 


(Q_A F_1,F_1)^{1/2}+(Q_A F_1,F_1)]. 


\end{gather*} 


 


Пусть величина $\lambda_1/(\lambda^2_2\alpha_1)\equiv\varepsilon$ мала, 


так чтобы было верно неравенство 


\begin{equation} 


\varepsilon+\sqrt{\varepsilon^2+\frac{\alpha_1}{4}\left(1+ 


\sqrt{\varepsilon}\sqrt{\alpha_1\lambda_1}\right)^2}\,\le 


1-\varepsilon_1, 


\end{equation} 


тогда будет иметь место априорная оценка 


$$ 


\varepsilon_1[(Ax,x)+\alpha_1(Q^{-1}_A B^*y,B^*y)]\le\Phi_3. 


$$ 


Для справедливости (3.17) можно потребовать выполнения условий 


$\alpha_1\le1$, $\alpha_1=2\sqrt{\alpha}$, $\varepsilon=1/9$, 


тогда $\varepsilon_1=0.009$, $p=9/\alpha_1$. 


\end{pf} 


 


Как видно из доказательства, неравенства (3.17), (3.8), (3.11), 


(3.12), (3.15) дают оценку $(D^{-1}L-\lambda I)^{-1}$ на границе 


$\Gamma$ области $\Pi$. Заметим, что 


оценка спектра типа (3.1) приведена в ([8], теорема 6). 


 


 


\section{Выбор параметров, оценка скорости сходимости} 


 


Заметим, что при $\alpha_1<4$ метод (1.2) при достаточно малых $\tau$ 


сходится. Это 


следует из п.\,2 доказательства теоремы 3. Используя представление 


степеней оператора пересчета $T$, для погрешности $z_n-z$ можно получить 


оценку 


$$ 


\|z_n-z\|\le C\max_{\lambda\in\Pi}|1-\tau\lambda|^n\equiv C q^n, 


$$ 


где норму можно составить из скалярных произведений, фигурирующих в 


условиях 


разрешимости (1.3), а $C$ определяется величиной интеграла 


$$ 


\oint\|(D^{-1}L-\lambda I)^{-1}\|\,|d\lambda|, 


$$ 


необходимые неравенства приведены в теореме 3. Для хорошего выбора 


параметров $\tau$, $\alpha$, $\alpha_1$ воспользуемся решением 


оптимизационной задачи 


из теоремы 1. Отметим, что в нашем случае всегда будем находиться 


в рамках условия $p\ge\lambda^+_1-\lambda^-_1$. Действительно, 


\begin{equation} 


p=9/\alpha_1,\ \ \lambda^+_1=(1+\alpha_1/\alpha), 


\ \ \lambda^-_1\ge0.5\min\{\A/4,\ m/(\alpha4)\}=R/2, 


\ \ \alpha_1=2\sqrt{\alpha},\ \ \alpha_1\le1. 


\end{equation} 


 


Качественно опишем рекомендации по выбору параметров. Если $\A\ll m$, то 


$\tau$, $\alpha$, $\alpha_1$ нужно брать около единицы, если $m\ll\A$, 


то $\tau\approx\alpha_1$ нужно брать 


малыми, а $\alpha$ --- еще меньшими. Отметим, что соотношение 


$\alpha_1=2\sqrt{\alpha}$ предельно 


точное в обоих случаях. Во всех вариантах скорость сходимости, грубо 


говоря, может быть в 100 раз хуже оптимальной. Дадим строгое описание 


результатов. 


 


Если взять параметр 


\begin{equation} 


\tau=\tau_-=(R/2+9/\alpha_1), 


\end{equation} 


то показатель скорости сходимости $q$ равен $q^2=(1+R\alpha_1/18)^{-1}$. 


Выберем $\alpha^+_1$ из условия максимизации 


$\alpha_1R=\min\{\alpha_1\A/4,\ m/\alpha_1\}$. 


Легко проверить, что 


\begin{equation} 


\alpha^+_1=1\ \ \text{при}\ \ 2\sqrt{m/\A}\ge1;\quad 


a^+_1=2\sqrt{m/\A}\ \ \text{при} \ \ 2\sqrt{m/\A}\le1. 


\end{equation} 


Тем самым доказана основная 


 


\begin{thm} 


Предположим, что верны условия разрешимости $(1.3)$ задачи $(1.1)$ и 


$z_n$ 


определены итерационным алгоритмом $(1.2)$. Пусть $0<\alpha_1\le1$, 


$\tau>0$, $\alpha_1=2\sqrt{\alpha}$, 


тогда для погрешности $z_n-z$ верна оценка 


$$ 


\|z_n-z\|\le c\rho^n(\tau)\|z_0-z\|,\quad n=1,2,\dotsc, 


$$ 


где $\rho(\tau)$ определяется соотношениями $(2.4)$, $(4.1)$. 


Если выбрать $\alpha_1=\alpha^+_1$, $\tau=\tau_-$, 


$\alpha_1=2\sqrt{\alpha}$ $($см. $(4.3)$, $(4.2))$, то $\|z_n-z\|\le c 


q^n_+\|z_0-z\|$, 


где 


$$ 


q^2_+=(1+\A/72)^{-1}\ \ \text{при} \ \ 2\sqrt{m/\A}\ge1; \quad 


q^2_+=(1+\sqrt{m/\A}\,\big/36)^{-1}\ \ \text{при} 


\ \ 2\sqrt{m/\A}\le1. 


$$ 


\end{thm} 


 


Насколько $q_+$ отличается от оптимального решения $q_0(\A,m)$ модельной 


задачи 


(1.4) в полном диапазоне изменения параметров $\A,m\in(0,1]^2$, 


позволяет судить 


теорема 2 (неравенства (2.5) или (2.6)). Наглядное представление этих 


оценок дает их асимптотика при малых $m$, $\A$ --- неравенства (1.7). 


 


 


\begin{thebibliography}{99} 


 


\bibitem{1} 


Fortin M., Glowinski R. {\em Augmented Lagrangian methods. 


Applications to the numerical solution of boundary value problems} // 


Stud. Math. and Appl. --


Amsterdam. -- 1983. -- \No\,15. -- 340 p. 


\bibitem{2} 


Чижонков Е.В. {\em К решению алгебраической системы типа Стокса при 


блочно-диагональном предобусловливании} // Журн. вычисл. 


матем. и матем. физ. -- 2001. -- Т.\,41. -- \No\,4. -- 


С.\,549--557. 


\bibitem{3} 


Queck W. {\em The convergence factor of preconditioned algorithms of 


the Arrow--Hurwicz type} // SIAM J. Numer. Anal. -- 1989. -- V.\,4. -- 


P.\,1016--1030. 


\bibitem{4} 


Silvester D., Wathen A. {\em Fast iterative solutions of stabilized 


Stokes systems}. II: {\em Using block diagonal preconditioners} // 


SIAM J. Numer. Anal. -- 1994. -- V.\,5. -- P.\,1352--1367. 


\bibitem{5} 


Bramble J.H., Pasciak J.E., Vassilev A.T. 


{\em Analysis of the inexact Uzawa algorithm for 


saddle point problems} // SIAM J. Numer. Anal. -- 1997. -- V.\,3. -- 


P.\,1072--1092. 


\bibitem{6} 


Чижонков Е.В. {\em Итерационные методы решения сеточных уравнений с 


седловым оператором}. Дис. $\dotsc$ докт. физ.-матем. наук. -- М.: МГУ, 


1999. -- 32 с. 


\bibitem{7} 


Астраханцев Г.П. {\em Об одном способе приближенного решения задачи 


Дирихле для бигармонического уравнения} // Журн. вычисл. матем. 


и матем. физ. -- 1977. -- Т.\,17. -- \No\,4. -- 


С.\,980--988. 


\bibitem{8} 


Астраханцев Г.П. {\em Анализ алгоритмов типа Эрроу--Гурвица} // 


Журн. вычисл. матем. и 


матем. физ. -- 2001. -- Т.\,41. -- \No\,1. -- С.\,17--28.


 


\end{thebibliography} 


 


\vskip 2cm 


 


\post{\it Санкт-Петербургский экономико-}{\it Поступила} 


\post{\it математический институт}{18.06.2002} 


\post{\it Российской Академии наук}{} 


 


\label{end} 


\end{document} 


